(1) 定义 \(M=m_1+m_2\)、质心 \(X=(m_1x_1+m_2x_2)/M\)、约化质量 \(\mu=m_1m_2/M\)。系统Lagrange函数为 \(L=M\dot X^2/2+\mu\dot\xi^2/2-k\xi^2/2\)，其中 \(x_2-x_1=\ell+\xi\)，无外力故质心自由运动。

(2) 正常坐标满足 \(\ddot X=0\)、\(\mu\ddot\xi+k\xi=0\)，\(\omega=\sqrt{k/\mu}\)。初态 \(\dot X(0)=m_1v_0/M\)、\(\xi(0)=0\)、\(\dot\xi(0)=-v_0\)，因此 \(X(t)=X(0)+(m_1v_0/M)t\)，\(\boxed{\xi(t)=-(v_0/\omega)\sin\omega t}\)。

(3) 由 \(v_1=\dot X-(m_2/M)\dot\xi\)、\(v_2=\dot X+(m_1/M)\dot\xi\)，代入 \(\dot\xi=-v_0\cos\omega t\) 得 \(v_1=m_1v_0/M+(m_2/M)v_0\cos\omega t\)、\(v_2=(m_1/M)v_0(1-\cos\omega t)\)。弹簧能 \(U=k v_0^2\sin^2(\omega t)/(2\omega^2)\)，和相对动能之和恒为 \(\mu v_0^2/2\)。

(4) 首次非零时刻恢复原长为 \(t=\pi/\omega\)，此时 \(\cos\omega t=-1\)，故 \(\boxed{v_1'=(m_1-m_2)v_0/(m_1+m_2),\quad v_2'=2m_1v_0/(m_1+m_2)}\)。弹簧能此时为零，总动量和能量均守恒，正是完全弹性一维碰撞的末态；当 \(m_1=m_2\) 时两球交换速度。
